% ch8_d (书页 362-375 / p377-p390)
%--C-TAIL--

% ==== tail：8.7 节（Our Universe）的收尾文字，承接书页 361 末的式 (8.175) 之后的段落，
% ==== 至本文件第一个 \section（8.8 暴胀）之前；由装配器移至 ch8_c 之后 ====

% ---- 书页 362 / p377 ----
事实上，只要相对论性自由度还是粒子物理标准模型中的那些，$g_{*}$ 与 $g_{*S}$ 就应当近似相等。一个十分粗略的估算指南是
\begin{equation*}
g_{*} \approx g_{*S} \sim
\left\{
\begin{array}{ll}
100 & T > 300~\mathrm{MeV} \\
10 & 300~\mathrm{MeV} > T > 1~\mathrm{MeV} \\
3 & T < 1~\mathrm{MeV}.
\end{array}
\right.
\tag{8.176}
\end{equation*}
正如我们马上就要讨论的，改变相对论性自由度有效数目的事件，是 $300~\mathrm{MeV}$ 处的 QCD 相变，以及 $1~\mathrm{MeV}$ 处电子/正电子对的湮灭。

有了这些背景，让我们来考虑宇宙从早期到今天的演化。作为开始，我们想象一个 Robertson--Walker 度规，其物质场处于温度 $1~\mathrm{TeV}=1000~\mathrm{GeV}$ 的热平衡中。高温等离子体是基本粒子（夸克、轻子、规范玻色子和 Higgs 玻色子）的复杂混合物。能量密度的主导形式将是相对论性粒子，因此早期宇宙是辐射主导的。它还非常接近平直，因为弗里德曼方程中的曲率项比物质密度和辐射密度演化得更慢。于是弗里德曼方程为
\begin{equation*}
\begin{aligned}
H^{2} &= \frac{8\pi G}{3}\rho_{\mathrm{R}} \\
&\approx 0.1\,g_{*}\,\frac{T^{4}}{\bar{m}_{\mathrm{P}}^{2}},
\end{aligned}
\tag{8.177}
\end{equation*}
其中约化 Planck 标度为 $\bar{m}_{\mathrm{P}}=(8\pi G)^{-1/2}\approx 10^{18}~\mathrm{GeV}$。如果辐射主导的阶段一直延伸到极早期，宇宙的年龄将近似为 $t\sim H^{-1}$，即
\begin{equation*}
t \sim \frac{\bar{m}_{\mathrm{P}}}{T^{2}}.
\tag{8.178}
\end{equation*}
用常规单位表示，这变为
\begin{equation*}
t \sim 10^{-6}\left(\frac{\mathrm{GeV}}{T}\right)^{2}~\mathrm{sec}.
\tag{8.179}
\end{equation*}

目前粒子加速器上的实验已经为大约 $100~\mathrm{GeV}$ 以下的物理图景提供了准确的描述，因此再外推一个数量级仍属于合理外推的范围。在更高的温度下我们就不太确定会发生什么了；在 $1~\mathrm{TeV}$ 和 Planck 标度之间也许没有任何非常有趣的事情，也可能这个区域充满了各式各样的意外。当然，也可以设想宇宙学在更低的温度下就会给出意外，尽管标准模型物理已被理解得很透彻；本节我们描述的是一个保守的方案，但一如既往，保持开放的头脑总是值得的。

% ---- 书页 363 / p378 ----
标准模型的一个关键特征是电弱部分的自发破缺对称性。在宇宙学中，这种对称性破缺发生在电弱相变处，温度 $T\sim 200~\mathrm{GeV}$。在此温度之上对称性尚未破缺，因此基本费米子（夸克和轻子）与弱相互作用规范玻色子都无质量；在此温度之下，我们便得到低能实验所熟知的那些质量模式。电弱相变预计不会给晚期宇宙留下任何可察觉的影响；一个可能的例外是重子生成（baryogenesis），下面讨论。

在这些温度下，量子色动力学（QCD）所描述的强相互作用并不那么强。在低能量/温度下，QCD 表现出“禁闭”（confinement）——夸克和胶子被束缚成重子和介子这样的复合粒子。但在 QCD 标度 $\Lambda_{\mathrm{QCD}}\sim 300~\mathrm{MeV}$ 之上，夸克和胶子是自由粒子。随着宇宙膨胀和冷却，强相互作用粒子被禁闭成束缚态，这正是 (8.176) 中所指出的相对论性自由度有效数目第一次下降的原因。QCD 相变预计不会在可观测宇宙上留下显著的印记。

正如强相互作用在高温下并不很强，弱相互作用也不像你可能以为的那么弱；它们仍然是弱的——在微扰论能够准确描述的意义上——但其发生得足够快，足以使中微子这类弱相互作用粒子保持在热平衡中。当 $T\sim 1~\mathrm{MeV}$ 时，情况不再如此。这也大致是电子和正电子变为非相对论性并发生湮灭的温度，相对论性自由度的有效数目因此减少，但这两件事互不相干。对 $1~\mathrm{MeV}$ 以下的温度，我们说弱相互作用已经“冻结”（frozen out）——相互作用速率降到宇宙的膨胀速率之下，相互作用发生得过于稀少，无法再使粒子保持平衡。冷暗物质粒子有可能就是在这个温度从等离子体中退耦的。更有把握的是，我们可以断定中子和质子不再相互转换。在这个温度下中子的平衡丰度约为质子丰度的 $1/6$（因为中子的质量稍大）。中子的寿命有限（$\tau_{n}=890~\mathrm{sec}$），只比这个时期宇宙的年龄 $t(1~\mathrm{MeV})\approx 1~\mathrm{sec}$ 略长，但它们已开始逐渐衰变为质子和轻子。然而不久之后，我们就到达略低于 $100~\mathrm{keV}$ 的温度，\textbf{大爆炸核合成}（Big-Bang Nucleosynthesis，BBN）开始了。

每个核子的核结合能典型地为 $1~\mathrm{MeV}$ 的量级，所以你也许会以为核合成应该更早发生；然而，每个核子对应的光子数太多，阻碍了核合成的发生，直到温度降到 $100~\mathrm{keV}$ 以下。在那一点上，中子/质子比约为 $1/7$。在所有轻核之中，核子驻留在 $^{4}\mathrm{He}$ 里在能量上是有利的，事实上大多数自由中子也确实被转化成了 $^{4}\mathrm{He}$；每两个中子和十四个质子，最终得到一个氦核和十二个质子。于是，按质量计约有 $25\%$ 的重子被转化为氦。此外，还有痕量的氘（每质子约 $10^{-5}$ 个氘核）、$^{3}\mathrm{He}$（同样约 $10^{-5}$）以及 $^{7}\mathrm{Li}$（约 $10^{-10}$）。

% ---- 书页 364 / p379 ----
当然，这些数字是预言，而它们确实被对轻元素原初丰度的观测所证实。（更重的元素并非在大爆炸中合成，而是需要晚期宇宙中的恒星过程。）我们对许许多多关键细节作了粗略带过，尤其是那些解释各种丰度如何依赖宇宙学参数的细节。例如，设想我们偏离标准模型，引入三种以上的轻中微子物种。这会通过式 (8.174) 在固定温度下增大辐射能量密度，转而又缩短了与给定温度相联系的时标（因为 $t\sim H^{-1}\propto\rho_{\mathrm{R}}^{-1/2}$）。核合成因此会稍早发生，导致中子丰度更高，从而 $^{4}\mathrm{He}$ 丰度也更大。对原初氦丰度的观测与标准模型的预言相符，给出了轻中微子数目接近三的第一个证据。类似地，与核合成相联系的所有温度和时标都依赖于重子/光子比；要与观测到的丰度相符，就要求每光子大约有 $5\times 10^{-10}$ 个重子，这正是重子密度参数之估计值 (8.164) 的由来，以及随之而来的对非重子暗物质的需要。

就我们当下的目的而言，原初核合成最深刻的特征，也许是它对温度与膨胀率之间的弗里德曼关系——从而对 Einstein 方程——的敏感依赖。BBN 的成功为 GR 提供了一个在远离我们日常经验的领域中的严格检验。Einstein 的理论当初被提出，主要是出于把引力与电磁学的洛伦兹对称性不变性相协调的需要；而这样一个理论竟能成功描述宇宙仅一秒钟大小时的增长，这实在是令人赞叹的成就。时至今日，BBN 仍是对各种替代引力理论最强有力的约束之一；特别地，它是我们拥有任何直接观测特征的最早时期。

核合成之后，我们得到一个由质子、电子和光子主导的等离子体，其中还有一些氦和其他核。也有暗物质存在，但假定到这个时期它已不与普通物质相互作用。下一个重要事件直到\textbf{复合}（recombination）才发生：电子与质子结合（它们与氦的结合稍早一些）。复合发生在温度 $T\approx 0.3~\mathrm{eV}$；此时宇宙是物质主导的。同样地，由于氢的结合能是 $13.6~\mathrm{eV}$，你也许会以为复合应该更早发生，但巨大的光子/重子比把它推迟了。复合至关重要的意义在于，它标志着宇宙变得透明的时刻。环境光子与自由电子强烈地相互作用，因此在复合之前光子的平均自由程非常短；而一旦电子和质子结合成中性氢，平均自由程实际上就变成了无穷大。这些环境光子今天作为宇宙微波背景被我们看到，它提供了宇宙在 $T\approx 0.3~\mathrm{eV}$ 时、即红移 $z\approx 1200$ 时的一幅快照。复合是一个颇为渐进的过程，因此对它发生于何时所作的任何指定都必然是近似的。

% ---- 书页 365 / p380 ----
复合之后，宇宙要走过一段被称为“黑暗时代”（dark ages）的漫长时期，其间星系通过引力不稳定性逐渐集结起来，但还没有可见的恒星来照亮宇宙。黑暗时代是一个神秘的时期；恒星和星系形成的过程高度复杂且非线性，无疑需要新类型的观测，我们才能对这一时代有充分的理解。

我们的故事现在已经把我们带到了今天，但还有几个缺失的环节需要回过头来补上。其一是宇宙中物质与反物质之间的不对称。宇宙中基本上全部可见物质似乎都由质子、中子和电子构成，而非它们的反粒子；如果遥远的星系主要由反物质构成，我们就应当能观测到物质/反物质区域边界处质子与反质子偶尔湮灭所产生的高能光子。虽然可以把这种不对称作为初始条件硬性加进去，但这总显得有些不能令人满意，大多数物理学家更愿意找到某种重子生成的动力学机制，让一个初始时物质/反物质对称的态演化成我们今天的宇宙。这样的破缺对称在粒子物理中很常见，事实上人们已经提出了许多重子生成机制（通常在电弱标度或更高的温度下）。然而，这些具体方案中没有一个已被证明足够令人信服而被采纳为标准方案。要理解重子不对称的起源，我们很可能需要对超越标准模型的物理有更好的理解。

另一个需要提及的缺失特征是，宇宙当然并非完全均匀和各向同性；宇宙中目前的大尺度结构似乎演化自极早期就存在的、幅度为 $\delta\rho/\rho\sim 10^{-5}$ 的绝热且近无标度微扰。这些微扰的绝热性和近无标度性的证据，来自对 CMB 和大尺度结构的联合观测。高度的各向同性和均匀性，以及对其的小偏离，在传统宇宙学中都只是作为神秘的初始条件被强加上去的。暴胀宇宙方案（inflationary-universe scenario）为两者提供了一个可能的动力学起源，我们现在就转向它。

\section{暴胀}

在对大爆炸模型的传统理解中，宇宙被认为在早期是辐射主导的，在晚期是物质主导的，而且——正如我们现在怀疑的——在非常晚近的时候才转变为真空主导。这幅图像在描述形形色色的观测数据方面取得了巨大成功；尽管如此，我们仍然可以问，产生这样一个宇宙的初始条件是否显得自然。这是那种人们在宇宙学中会问、在其他科学中却不会问的问题。通常，作为物理学家，我们寻找自然定律，并想象我们可以自由地指定初始条件，然后问它们在这些定律下如何演化。但宇宙似乎只有一组初始条件，因此，想知道它们究竟是相对平常的还是经过精细调节的，似乎是合情合理的。

% ---- 书页 366 / p381 ----
在传统图像中，早期宇宙确实被精细调节到了难以置信的精度。特别是，我们宇宙的两个特征看上去非常不寻常：它的空间平直性，以及它高度的各向同性和均匀性。有可能这就是我们恰好被给予的宇宙，追问不同初始条件的可能性有多大毫无意义；另一种可能是，如果存在某种动力学机制，能把宽广谱系的初始条件朝平直性和均匀性/各向同性的方向演化，那么这些条件就比乍看之下更有可能出现。暴胀宇宙方案提供了这样一种机制（而且还不止如此），它已成为现代宇宙学的核心组织原理，即使我们还远不能证明它的正确性。

在描述暴胀之前，让我们先描述它声称解决的两个不自然性问题：与均匀性/各向同性相联系的\textbf{平坦性问题}（flatness problem）和\textbf{视界问题}（horizon problem）。平坦性问题来自考虑一个有物质和辐射但没有真空能的宇宙中的弗里德曼方程；为方便后文，我们用约化 Planck 质量 $\bar{m}_{\mathrm{P}}=(8\pi G)^{-1/2}$ 把它写成
\begin{equation*}
H^{2} = \frac{1}{3\bar{m}_{\mathrm{P}}^{2}}\left(\rho_{\mathrm{M}}+\rho_{\mathrm{R}}\right) - \frac{\kappa}{a^{2}}.
\tag{8.180}
\end{equation*}
曲率项 $-\kappa/a^{2}$ 与 $a^{-2}$ 成正比（这显而易见），而能量密度项随尺度因子的增大衰减得更快，$\rho_{\mathrm{M}}\propto a^{-3}$、$\rho_{\mathrm{R}}\propto a^{-4}$。这就引出了一个问题：考虑到自 Planck 时期以来 $a$ 已经增大了大约 $10^{30}$ 倍，为什么比值 $(\kappa a^{-2})/(\rho/3\bar{m}_{\mathrm{P}}^{2})$ 没有远远大于 1？换句话说，在物质/辐射主导的宇宙中，$\Omega=1$ 这个点是一个排斥性的不动点——对这个值的任何偏离都会随时间增长，那么为什么我们今天观测到的是 $\Omega\approx 1$？

视界问题源于 FRW 宇宙学中粒子视界的存在，如图 8.7 所示。视界之所以存在，是因为自大爆炸奇点以来只有有限的时间，因而在宇宙的年龄之内光子只能走过有限的距离，正如我们在第 2 章中简要讨论过的。考虑在平直宇宙中沿径向轨迹运动的一个光子（推广到非平直宇宙是直截了当的）。径向类光路径满足
\begin{equation*}
0 = \mathrm{d}s^{2} = -\mathrm{d}t^{2} + a^{2}\mathrm{d}r^{2},
\tag{8.181}
\end{equation*}
因此这样的一个光子在 $t_{1}$ 与 $t_{2}$ 两时刻之间走过的共动（坐标）距离为
\begin{equation*}
\Delta r = \int_{t_{1}}^{t_{2}}\frac{\mathrm{d}t}{a(t)}.
\tag{8.182}
\end{equation*}
要得到任一时刻 $t$ 处观者所测得的物理距离，只需乘上 $a(t)$。为简单起见，让我们想象自己处在一个物质主导的
% ---- 书页 367 / p382 ----
宇宙，对它而言
\begin{equation*}
a = \left(\frac{t}{t_{0}}\right)^{2/3}.
\tag{8.183}
\end{equation*}
%FIG{8.7}: {从一个由大爆炸奇点膨胀而来的宇宙中的过去光锥，展示了宇宙学中的粒子视界。处于复合时刻的一些点——今天作为天空中相反两侧的宇宙微波背景被我们观测到——其过去光锥互不重叠（在传统宇宙学中）；没有任何因果信号能够影响它们，使它们具有相同的温度。}
记住 $a_{0}=1$。因此 Hubble 参量为
\begin{equation*}
\begin{aligned}
H &= \frac{2}{3}t^{-1} \\
&= a^{-3/2}H_{0}.
\end{aligned}
\tag{8.184}
\end{equation*}
于是该光子走过的共动距离为
\begin{equation*}
\Delta r = 2H_{0}^{-1}\left(\sqrt{a_{2}}-\sqrt{a_{1}}\right).
\tag{8.185}
\end{equation*}
在尺度因子取任一固定值 $a=a_{*}$ 时的共动视界尺度，是光子自大爆炸以来走过的距离，
\begin{equation*}
r_{\mathrm{hor}}(a_{*}) = 2H_{0}^{-1}\sqrt{a_{*}}.
\tag{8.186}
\end{equation*}
而在 $a_{*}$ 处的空间超曲面上测得的物理视界尺寸，因此就简单地是
\begin{equation*}
d_{\mathrm{hor}}(a_{*}) = a_{*}r_{\mathrm{hor}}(a_{*}) = 2H_{*}^{-1}.
\tag{8.187}
\end{equation*}
事实上，对任何包含物质和辐射混合物的近平直宇宙，在任一给定时期我们都有
\begin{equation*}
d_{\mathrm{hor}}(a_{*}) \sim H_{*}^{-1} = d_{H}(a_{*}),
\tag{8.188}
\end{equation*}

% ---- 书页 368 / p383 ----
其中 Hubble 距离 $d_{H}$ 已在式 (8.71) 中引入。这种近似相等带来一种强烈的诱惑，即把“视界距离”和“Hubble 距离”当作同义词使用；这种诱惑应当抵制，因为暴胀可以使前者远大于后者，我们很快就会证明这一点。

视界问题简单说来就是：CMB 在很高的精度上是各向同性的，尽管最后散射面上相距遥远的各点彼此完全处于对方的视界之外。当我们观看 CMB 时，我们观测的是尺度因子 $a_{\mathrm{CMB}}\approx 1/1200$ 处的宇宙；由式 (8.185)，CMB 上一点与地球上的观者之间的共动距离为
\begin{equation*}
\begin{aligned}
\Delta r &= 2H_{0}^{-1}\left(1-\sqrt{a_{\mathrm{CMB}}}\right) \\
&\approx 2H_{0}^{-1}.
\end{aligned}
\tag{8.189}
\end{equation*}
然而，这样一个点所对应的共动视界距离为
\begin{equation*}
\begin{aligned}
r_{\mathrm{hor}}(a_{\mathrm{CMB}}) &= 2H_{0}^{-1}\sqrt{a_{\mathrm{CMB}}} \\
&\approx 6\times 10^{-2}H_{0}^{-1}.
\end{aligned}
\tag{8.190}
\end{equation*}
因此，如果我们观测 CMB 中相距遥远的两个部分，它们的视界将互不重叠；CMB 天空中不同的斑块在复合时刻是因果不连通的。然而，观测却表明它们具有相同的温度，且精度很高。于是问题就来了：既然它们从未处于因果接触之中，它们是如何预先知道要以正确的方式协调各自的演化的？我们必须设法修正传统 FRW 宇宙学的因果结构。

让我们考虑这样修正传统图像：设定一段\textbf{暴胀}（inflation）时期，即极早期宇宙中一段加速膨胀（$\ddot{a}>0$）的时期，由物质或辐射以外的某种成分驱动，这种成分随宇宙膨胀而红移得很慢。这样一来，平坦性问题和视界问题便可以同时解决。为简单起见，考虑暴胀由恒定的真空能驱动、从而导致指数膨胀的情形。于是，在真空主导的时期，$\rho/3\bar{m}_{\mathrm{P}}^{2}\propto a^{0}$ 相对于 $-\kappa/a^{2}$ 迅速增大，宇宙随时间变得越来越平（$\Omega$ 被驱动到 1）。如果这个过程持续了足够长的时间，然后真空能才转化为物质和辐射，那么密度参数将足够接近 1，以至于进入当前时代之前它还来不及发生可察觉的改变。与此同时，视界问题可以追溯到这样一个事实：任意两个共动物体之间的物理距离随尺度因子增长，而物质主导或辐射主导宇宙中的物理视界尺寸增长得更快，$d_{\mathrm{hor}}\sim a^{n/2}H_{0}^{-1}$。这同样可以由早期一段指数膨胀时期来解决：真实的视界尺寸会增长到一个惊人的数量，以至于我们今天的视界实际上远大于“视界等于 Hubble 半径 $H_{0}^{-1}$”这个天真的估计。

事实上，真正的指数膨胀并非必需；对任何加速膨胀，空间曲率相对于能量密度都会减小，
% ---- 书页 369 / p384 ----
而视界距离将迅速增长。通常我们要求这个加速时期持续 60 个或更多的 $e$ 指数倍（e-fold，$e$ 指数倍数为 $N=\Delta\ln a$），这正是解决视界问题所需要的。很容易做过头，暴胀一般会使今天的宇宙在空间上平直得难以置信。

现在让我们考虑如何在早期宇宙中得到一个暴胀阶段。最直接的办法是利用一个标量场的势所提供的真空能，这个场就是\textbf{暴胀子}（inflaton）。想象一个由空间均匀的标量（场）的能量主导的宇宙。相关的运动方程正是我们在 8.7 节中讨论动力学暗能量时的那些方程；唯一的区别是暴胀的能量标度要高得多。我们有 RW 度规中标量场的运动方程，
\begin{equation*}
\ddot{\phi} + 3H\dot{\phi} + V'(\phi) = 0,
\tag{8.191}
\end{equation*}
以及弗里德曼方程
\begin{equation*}
H^{2} = \frac{1}{3\bar{m}_{\mathrm{P}}^{2}}\left(\frac{1}{2}\dot{\phi}^{2} + V(\phi)\right).
\tag{8.192}
\end{equation*}
我们忽略了曲率项，因为暴胀反正会把宇宙拉平。如果场的演化足够平缓，使势能得以主导动能，并且 $\phi$ 的二阶导数足够小，足以让这种状态维持一段充分长的时间，暴胀就能发生。因此，我们希望
\begin{equation*}
\begin{aligned}
\dot{\phi}^{2} &\ll V(\phi), \\
|\ddot{\phi}| &\ll |3H\dot{\phi}|,\ |V'|.
\end{aligned}
\tag{8.193}
\end{equation*}
满足这些条件要求两个无量纲量足够小，它们就是所谓的\textbf{缓慢滚动参数}（slow-roll parameters）：
\begin{equation*}
\begin{aligned}
\epsilon &= \frac{1}{2}\bar{m}_{\mathrm{P}}^{2}\left(\frac{V'}{V}\right)^{2}, \\
\eta &= \bar{m}_{\mathrm{P}}^{2}\left(\frac{V''}{V}\right).
\end{aligned}
\tag{8.194}
\end{equation*}
注意 $\epsilon\ge 0$，而 $\eta$ 可正可负。还要注意，这些定义并非人人通用；有些人喜欢用 Hubble 参量而不是势来定义它们。我们的选择描述的是场是否有机会缓慢滚动一阵子；而用 Hubble 参量的描述给出的则是场实际上是否正在缓慢滚动。当这两个量都很小时，我们就可以有一个长久的暴胀阶段。然而它们并不充分；无论势是什么样子，我们总可以选取初始条件使 $|\dot{\phi}|$ 大到缓慢滚动永不适用。不过，只要缓慢滚动参数很小，大多数初始条件都会被吸引到暴胀阶段去。

% ---- 书页 370 / p385 ----
要发明满足缓慢滚动条件的势并不难。考虑也许是最简单的例子\footnote{我们沿用 A.R.~Liddle, \emph{An Introduction to Cosmological Inflation}（\texttt{http://arxiv.org/astro-ph/9901124}）中的讲述。}，$V(\phi)=\frac{1}{2}m^{2}\phi^{2}$。在这种情形下
\begin{equation*}
\epsilon = \eta = \frac{2\bar{m}_{\mathrm{P}}^{2}}{\phi^{2}}.
\tag{8.195}
\end{equation*}
显然，只要 $\phi$ 足够大，我们就能让缓慢滚动参数要多小有多小。然而，我们有一个约束：能量密度不应高达 Planck 标度，这样我们的经典分析才有意义；这意味着 $\phi\ll\bar{m}_{\mathrm{P}}^{2}/m$。如果我们把场从值 $\phi_{i}$ 开始，那么在暴胀结束之前（也就是缓慢滚动参数变成 1 的量级之前）的 $e$ 指数倍数将是
\begin{equation*}
\begin{aligned}
N &= \int_{t_{i}}^{t_{e}} H\,\mathrm{d}t \\
&\approx -\bar{m}_{\mathrm{P}}^{-2}\int_{\phi_{i}}^{\phi_{e}}\frac{V}{V'}\,\mathrm{d}\phi \\
&\approx \frac{\phi_{i}^{2}}{4\bar{m}_{\mathrm{P}}^{2}} - \frac{1}{2}.
\end{aligned}
\tag{8.196}
\end{equation*}
第一个等式总是成立，第二个等式用了缓慢滚动近似，第三个等式是这个具体模型的结果。要得到 60 个 $e$ 指数倍，我们需要 $\phi_{i}>16\bar{m}_{\mathrm{P}}$。再加上能量密度的上限，我们发现质量参数存在一个上限：$m\ll\bar{m}_{\mathrm{P}}/16$。事实上，观测到的密度涨落的大小对 $m$ 给出了更严格的上限，我们将在下面讨论。但 $m$ 没有下限，因此，只有当我们愿意设定 $m\ll\bar{m}_{\mathrm{P}}$ 这样的巨大层级差，或者等价地一个很小的无量纲数 $m/\bar{m}_{\mathrm{P}}$ 时，才容易得到合适的暴胀势。用 $\lambda\phi^{4}$ 势做同样的练习也会得出类似的结论：$\lambda$ 必须相当小；我们常说暴胀子必须是弱耦合的。当然，从某种意义上说这是在作弊，因为当场值 $\phi>\bar{m}_{\mathrm{P}}$ 时，我们应当预期有效势中会出现形如 $\bar{m}_{\mathrm{P}}^{4-n}\phi^{n}$（$n>4$）的附加项并变得重要。所以在一个\emph{现实}的模型中，要得到合适的势可能相当困难。

暴胀终会在某个时刻结束，暴胀子势中的能量将转化为物质和辐射的热化气体，这个过程被称为“再热”（reheating）。恰当地理解再热过程至关重要，因为它控制着我们宇宙中种种我们想要或不想要的遗留物（relic）的产生。例如，暴胀的一个重要好处是，它能够把早期宇宙中可能产生、但今天并未观测到的各种遗留物“暴胀掉”（inflate away）。一个经典的例子出现在粒子物理大统一理论的语境中：大统一理论一般预言存在超重的磁单极子，其丰度比观测容许的还要多出许多个数量级。

% ---- 书页 371 / p386 ----
历史上，磁单极子问题是 Guth 发明暴胀的首要动机；平坦性问题和视界问题的解决当时被视为额外的红利。暴胀可以把磁单极子的丰度稀释到合适的程度，但如果宇宙再热到超过大统一相变的温度，它们又会重新产生；幸运的是，对大多数模型而言这并不是一个严格的约束。类似的考虑也适用于其他不想要的遗留物；在超对称模型中，引力微子（gravitino，引力子的超对称伙伴）的丰度引起一个尤其令人头疼的问题。与此同时，又必须再热到足够高的温度，以便某种重子生成方案得以进行。对于在某个粒子物理模型内实现暴胀的任何具体方案，关键的检查是：不想要的遗留物被清除，而想要的遗留物（例如重子）得以保留。

暴胀方案的一个关键要素是密度微扰的产生，它们可能是我们今天观测到的 CMB 温度各向异性和星系大尺度结构的起源。暴胀产生密度微扰背后的想法相当直截了当。暴胀会把任何环境粒子密度迅速衰减到零，只留下真空。但是加速宇宙中的真空态具有非零的温度，即 Gibbons--Hawking 温度，它类似于黑洞的 Hawking 温度。我们无法在本书中深入探讨这个主题；这里只概述其基本结果。

对于一个由势能 $V$ 主导的宇宙，Gibbons--Hawking 温度为
\begin{equation*}
T_{\mathrm{GH}} = \frac{H}{2\pi} \sim \frac{V^{1/2}}{\bar{m}_{\mathrm{P}}}.
\tag{8.197}
\end{equation*}
与这个温度相对应的，是暴胀子场 $\phi$ 在每个波数 $k$ 处的涨落，其大小为
\begin{equation*}
|\Delta\phi|_{k} = T_{\mathrm{GH}}.
\tag{8.198}
\end{equation*}
由于势按假设近乎平坦，$\phi$ 的涨落会导致能量密度的小涨落，
\begin{equation*}
\delta\rho = V'(\phi)\delta\phi.
\tag{8.199}
\end{equation*}
因此，暴胀在所有尺度上都产生密度微扰。微扰的幅度在每个波数处都几乎相等，但由于暴胀子滚动时 $V$ 的逐渐变化，会有微小的偏离。描述这些微扰是一个繁乱的课题，涉及数不清的各种记号。一个合理的出发点是方均根（root-mean-square，RMS）密度涨落，
\begin{equation*}
\left.\frac{\delta\rho}{\rho}\right|_{\mathrm{rms}} = \sqrt{\left\langle\left(\frac{\delta\rho}{\rho}\right)^{2}\right\rangle},
\tag{8.200}
\end{equation*}

% ---- 书页 372 / p387 ----
其中尖括号表示对空间位置的平均。对于统计各向同性的微扰（期望幅度与方向无关），稍稍做一点傅里叶分析便能让我们写出
\begin{equation*}
\left(\left.\frac{\delta\rho}{\rho}\right|_{\mathrm{rms}}\right)^{2} = \int \Delta^{2}(k)\,\mathrm{d}(\ln k),
\tag{8.201}
\end{equation*}
其中我们引入了无量纲功率谱，
\begin{equation*}
\Delta^{2}(k) \equiv \frac{k^{3}|\delta_{k}|^{2}}{2\pi^{2}},
\tag{8.202}
\end{equation*}
而 $\delta_{k}$ 是分数密度微扰（fractional density perturbation）的傅里叶变换的期望值，
\begin{equation*}
\delta_{\mathbf{k}} = \frac{1}{(2\pi)^{3/2}}\int e^{-i\mathbf{k}\cdot\mathbf{x}}\frac{\delta\rho}{\rho}\,\mathrm{d}^{3}x,
\tag{8.203}
\end{equation*}
这里我们已假设它是各向同性的。无量纲功率谱是时间的函数，因为每个模式的幅度都在演化；表达任何具体模型的预言时，最常用的做法是取模式的物理波长 $\lambda=a/k$ 等于 Hubble 半径 $H^{-1}$ 的那一时刻的微扰幅度，
\begin{equation*}
A_{S}^{2}(k) \equiv \left.\Delta^{2}(k)\right|_{k=aH}.
\tag{8.204}
\end{equation*}
这样，$A_{S}(k)$ 量度的是不同模式在不同时刻的幅度。对于由缓慢滚动标量场驱动的暴胀，$A_{S}(k)$ 与势通过下式相联系：
\begin{equation*}
A_{S}^{2}(k) \sim \left.\frac{V^{3}}{\bar{m}_{\mathrm{P}}^{6}(V')^{2}}\right|_{k=aH} \sim \left.\frac{V}{\bar{m}_{\mathrm{P}}^{4}\epsilon}\right|_{k=aH}.
\tag{8.205}
\end{equation*}
我们有意略去了无量纲的数值因子——它们在不同文献之间差别很大——以便突出对势的依赖。

谱带着下标“S”，是因为它描述的是度规中的标量涨落。这些涨落与能量-动量分布绑定在一起，暴胀产生的密度涨落是绝热的——所有物种的密度涨落都是相互关联的。涨落还是高斯的，其含义是：描述不同尺度上涨落的那些傅里叶模式的相位互不相关。暴胀微扰的这些方面——近无标度的绝热密度涨落谱、高斯分布——都与当前对 CMB 和大尺度结构的观测一致，而今后几年预定要收集的新数据应会大大提高这些检验的精度。

% ---- 书页 373 / p388 ----
在暴胀期间被激发的不仅是近无质量的暴胀子，还有任何其他近无质量的粒子。另一个重要的例子是引力子，它对应于度规中的张量微扰（引力场的传播激发）。张量涨落具有谱
\begin{equation*}
A_{T}^{2}(k) \sim \left.\frac{V}{\bar{m}_{\mathrm{P}}^{4}}\right|_{k=aH}.
\tag{8.206}
\end{equation*}
重要的是，张量幅度只依赖于势本身，而不依赖于它的导数；因此，对张量微扰的观测将直接给出关于暴胀能量标度的信息。

为了理解观测，用可观测的量来参数化微扰谱是有用的。于是我们写出
\begin{equation*}
A_{S}^{2}(k) \propto k^{n_{S}-1}
\tag{8.207}
\end{equation*}
以及
\begin{equation*}
A_{T}^{2}(k) \propto k^{n_{T}},
\tag{8.208}
\end{equation*}
其中 $n_{S}$ 和 $n_{T}$ 是谱指数（spectral index）。它们与势的缓慢滚动参数的关系为
\begin{equation*}
n_{S} = 1 - 6\epsilon + 2\eta
\tag{8.209}
\end{equation*}
以及
\begin{equation*}
n_{T} = -2\epsilon.
\tag{8.210}
\end{equation*}
在我们所考虑的这类模型（由单个缓慢滚动的标量场驱动）中，存在一个把标量模式和张量模式的幅度与谱指数联系起来的自洽关系。它可以用一种与约定无关的方式表达为可观测量之间的关系，即不同微扰所导致的温度涨落 $\Delta T$ 之间的关系：
\begin{equation*}
\frac{(\Delta T/T)_{T}^{2}}{(\Delta T/T)_{S}^{2}} = -7n_{T}.
\tag{8.211}
\end{equation*}

张量微扰的存在是暴胀的一个关键预言，原则上可以通过对 CMB 偏振的观测加以检验。偏振也可以由通常的密度涨落诱导产生，其机制是非均匀等离子体中 Thompson 散射截面的各向异性。幸运的是，我们可以设想把天空中的偏振矢量场分解为无旋部分（$E$ 模式）和有旋部分（$B$ 模式）；标量微扰导致 $E$ 模式偏振，而张量微扰导致 $B$ 模式（除去复合之后的宇宙中某些不可避免的演化加工）。CMB 偏振已经被探测到了；未来的挑战将在于把标量和张量的贡献分离出来，以检验简单暴胀模型的预言 (8.211)。

% ---- 书页 374 / p389 ----
当然，这不仅要求探测到张量诱导的偏振，还要求以一定的精度测量它的谱指数。

我们对微扰幅度已有的知识，已经给了我们关于暴胀能量标度的重要信息。张量微扰只依赖 $V$ 本身，不依赖它的导数；如果 COBE 观测到的 CMB 各向异性源自张量涨落（有可能，尽管不太可能），我们可以立刻导出 $V_{\mathrm{inflation}}\sim(10^{16}~\mathrm{GeV})^{4}$。这里，受到约束的 $V$ 值是产生观测涨落的那一时刻的值，即暴胀结束前 60 个 $e$ 指数倍时的值。这极其让人联想到大统一标度，十分令人鼓舞。即使在更可能的情形下——CMB 中观测到的微扰本质上是标量的——我们仍然可以写出
\begin{equation*}
V_{\mathrm{inflation}}^{1/4} \sim \epsilon^{1/4}10^{16}~\mathrm{GeV},
\tag{8.212}
\end{equation*}
其中 $\epsilon$ 是式 (8.194) 定义的缓慢滚动参数。虽然我们预期 $\epsilon$ 很小，但指数中的 $1/4$ 意味着对 $\epsilon$ 的依赖相当弱；除非这个参数小得出奇，否则很有可能 $V_{\mathrm{inflation}}^{1/4}\sim 10^{15}$--$10^{16}~\mathrm{GeV}$。我们能对如此惊人的能量标度获得这样的信息，这一事实本身就足以令人惊叹。

\section{习题}

\textbf{1.}\quad 考虑一个 $(N+n+1)$ 维时空，坐标为 $\{t,\,x^{I},\,y^{i}\}$，其中 $I$ 从 1 取到 $N$，$i$ 从 1 取到 $n$。设度规为
\begin{equation*}
\mathrm{d}s^{2} = -\mathrm{d}t^{2} + a^{2}(t)\delta_{IJ}\,\mathrm{d}x^{I}\mathrm{d}x^{J} + b^{2}(t)\gamma_{ij}(y)\,\mathrm{d}y^{i}\mathrm{d}y^{j},
\tag{8.213}
\end{equation*}
其中 $\delta_{IJ}$ 是通常的 Kronecker delta，而 $\gamma_{ij}(y)$ 是 $n$ 维最大对称空间流形上的度规。设想我们把度规 $\gamma$ 归一化，使得曲率参数
\begin{equation*}
k = \frac{R(\gamma)}{n(n-1)}
\tag{8.214}
\end{equation*}
等于 $+1$、$0$ 或 $-1$，其中 $R(\gamma)$ 是度规 $\gamma_{ij}$ 所对应的 Ricci 标量。

\textbf{(a)}\quad 计算这个度规的 Ricci 张量。

\textbf{(b)}\quad 用能量密度 $\rho$ 以及 $x^{I}$ 和 $y^{i}$ 方向上的压强 $p^{(N)}$、$p^{(n)}$ 定义一个能量-动量张量：
\begin{gather*}
T_{00} = \rho \tag{8.215}\\
T_{IJ} = a^{2}p^{(N)}\delta_{IJ} \tag{8.216}\\
T_{ij} = b^{2}p^{(n)}\gamma_{ij}. \tag{8.217}
\end{gather*}
把这个度规和 $T_{\mu\nu}$ 代入 Einstein 方程，推导 $a$ 和 $b$ 所满足的类弗里德曼方程（共三个独立方程）。

% ---- 书页 375 / p390 ----
\textbf{(c)}\quad 推导静态解（即 $\dot{a}=\dot{b}=\ddot{a}=\ddot{b}=0$ 的解）处能量密度和两个压强所满足的方程，用 $k$、$n$ 和 $N$ 表示。利用它们推导物态方程参数 $w^{(N)}=p^{(N)}/\rho$ 和 $w^{(n)}=p^{(n)}/\rho$ 在静态解处成立的表达式。

\textbf{2.}\quad 考虑 de Sitter 空间，采用度规取如下形式的坐标：
\begin{equation*}
\mathrm{d}s^{2} = -\mathrm{d}t^{2} + e^{2Ht}\left[\mathrm{d}x^{2}+\mathrm{d}y^{2}+\mathrm{d}z^{2}\right].
\tag{8.218}
\end{equation*}
对非共动观者（$x^{i}$ 不是常数）求解测地线方程，求出作为 $t$ 的函数的仿射参量。证明测地线在有限的仿射参量内到达 $t=-\infty$，从而证明这些坐标不能覆盖整个流形。

\textbf{3.}\quad 在附录 F 中我们讨论了 Raychaudhuri 方程。证明：把它应用到 Robertson--Walker 宇宙学上，Raychaudhuri 方程等价于第二弗里德曼方程 (8.68)。

\textbf{4.}\quad 考虑最佳拟合宇宙，其密度参数为 $\Omega_{\mathrm{R}0}=10^{-4}$、$\Omega_{\mathrm{M}0}=0.3$、$\Omega_{\Lambda0}=0.7$。在对数标度上画出三个 $\Omega_{i}$ 作为尺度因子 $a$ 的函数的图像，$a$ 从 $10^{-35}$ 取到 $10^{35}$。标出 Planck 时间、核合成和今天。

\textbf{5.}\quad 在平直时空中，物理尺寸固定的物体离得越远，所张的角就越小；在膨胀宇宙中未必如此。考虑红移 $z$ 处物理尺寸为 $L$ 的物体的角大小 $\theta(z)$。在物质主导的平直宇宙中，$\theta(z)/L$ 在什么红移处取极小？如果所有星系的跨度都至少有 $10~\mathrm{kpc}$（并且历来如此），这样一个宇宙中星系的最小角大小是多少？把你的结果既用 $H_{0}$ 表示，也代入 $H_{0}=70~\mathrm{km/s/Mpc}$ 给出数值。

\textbf{6.}\quad 在宇宙学中，我们倾向于把非相对论性粒子理想化为温度 $T$ 和压强 $p$ 都为零。实际上，随机运动会使它们具有一定的温度和压强，满足 $p\propto T\rho$。

\textbf{(a)}\quad 由有质量粒子组成的气体，其压强作为尺度因子的函数会如何衰减？

\textbf{(b)}\quad 设中微子具有质量 $m_{\nu}=0.1~\mathrm{eV}$，当前温度为 $T_{\nu 0}=2~\mathrm{K}$。中微子大约在什么红移从相对论性变为非相对论性？

\textbf{7.}\quad 设想宇宙在 Planck 时刻以均分状态（equipartition）出发（这样物质和辐射中的能量密度为 Planck 密度的量级，时间与空间的曲率半径为 Planck 长度的量级）。忽略任何空间不均匀性，计算一个正曲率宇宙能维持多久，以及一个负曲率宇宙在温度降到 $3~\mathrm{K}$ 时年龄多大。一个平直宇宙在温度降到 $3~\mathrm{K}$ 时年龄多大？一个平直宇宙到膨胀速率降至 $H_{0}=70~\mathrm{km\,s^{-1}\,Mpc^{-1}}$ 时年龄多大？
